\bookchapter{الفصل 22: كيف يعمل الذكاء الاصطناعي في الهاتف}

\chapterintro{مرحبًا بك في أروع “عرض سحري” على الهاتف، وبطل القصة هو: الذكاء الاصطناعي. الصور الجميلة “بضغطة واحدة”، والمساعد الصوتي الذي يجيب فورًا، ولوحة المفاتيح التي تفهمك كلما استخدمتها أكثر، كلها يقف خلفها هذا الدماغ الخفي.}

\booksection{أولًا: كاميرا الذكاء الاصطناعي — “غرفة التحميض” الخاصة بك انتقلت إلى الهاتف}

التصوير الحسابي بالذكاء الاصطناعي اليوم أشبه بإدخال فريق محترف من مصحح الصور والمصورين مباشرة في زر الغالق. قبل أن تضغط على الغالق، يكون الذكاء الاصطناعي قد بدأ العمل بالفعل: يمسح المشهد ويحدد “هذه سماء زرقاء” و“هذا وجه” و“هذه لقطة ليلية”.

\booksection{ثانيًا: المساعد الصوتي — كيف يفهمك “سكرتير الصوت” الخاص بك؟}

\begin{enumerate}
  \item \textbf{السماع بوضوح (التعرف على الكلام)}: تحويل الموجات الصوتية إلى نص.
  \item \textbf{الفهم (فهم اللغة الطبيعية)}: فهم القصد واستخراج الكيانات الأساسية.
  \item \textbf{التنفيذ والإجابة}: استدعاء التطبيق أو الخدمة المناسبة، ثم "التحدث" إليك عبر تركيب الكلام.
\end{enumerate}

\booksection{ثالثًا: الذكاء الاصطناعي على الجهاز مقابل الذكاء الاصطناعي السحابي}

يبقي الذكاء الاصطناعي على الجهاز الحسابات داخل الهاتف، فهو سريع وأكثر خصوصية ويعمل دون اتصال. أما الذكاء الاصطناعي السحابي فيتطلب اتصالًا بالإنترنت، لكنه يستطيع معالجة مهام أكثر تعقيدًا.

\booksection{رابعًا: النماذج الكبيرة للذكاء الاصطناعي والهاتف — مساعدك الذكي المحمول}

في 2025-2026، انفجار نماذج AI الكبيرة (مثل DeepSeek R1 وChatGPT-5 وERNIE 5.0 وQwen 3 وغيرها) غيّر عميقًا طريقة استخدام الهاتف. لم يعد الهاتف مجرد أداة لتنفيذ التعليمات، بل شريكاً ذكياً يفهم ويستنتج ويبدع. رقائق AI الجانبية (قوة حوسبة NPU تتجاوز 80 TOPS) جعلت استنتاج النماذج الكبيرة يعمل محلياً على الهاتف لأول مرة، يمكن استخدام AI دون اتصال بالإنترنت.

\booksubsection{ما هي النماذج الكبيرة للذكاء الاصطناعي؟}

النماذج الكبيرة للذكاء الاصطناعي (Large Language Model، LLM) هي أنظمة ذكاء اصطناعي مُدرَّبة بكميات هائلة من البيانات النصية، قادرة على فهم اللغة الطبيعية وتوليدها. يمكنك تخيلها كـ “طالب خارق” قرأ كل النصوص تقريبًا على الإنترنت، قادر على الإجابة عن الأسئلة وكتابة المقالات والترجمة والبرمجة وتحليل البيانات.

\booksubsection{تطبيقات النماذج الكبيرة للذكاء الاصطناعي على الهاتف:}

\begin{itemize}
  \item \textbf{مساعد المحادثة}: تطبيقات مثل DeepSeek وChatGPT وDoubao وKimi، يمكنك من خلالها طرح الأسئلة والترجمة وكتابة النصوص وتلخيص المقالات في أي وقت. إنها أكثر ذكاءً من محركات البحث التقليدية — فهي لا تعطيك روابط فحسب، بل تمنحك الإجابات مباشرة.
  \item \textbf{الكتابة والإبداع}: دع الذكاء الاصطناعي يكتب لك رسائل البريد الإلكتروني ويلخص محاضر الاجتماعات ويحسّن المقالات ويولد نصوصًا إبداعية.
  \item \textbf{توليد الصور وتحريرها}: يزيل الذكاء الاصطناعي بضغطة واحدة المارة في الصور، ويوسع حدود الصور، ويحول الأنماط، بل ويولد صورًا بناءً على الأوصاف النصية.
  \item \textbf{الترجمة الفورية}: ترجمة المكالمات الفورية، وترجمة المحادثات وجهًا لوجه، وترجمة الصور، فلم تعد هناك حواجز لغوية أثناء السفر للخارج.
  \item \textbf{المساعدة في التعلم}: يستطيع الذكاء الاصطناعي شرح المفاهيم وطرح أسئلة التدريب وتصحيح الواجبات، وهو بمثابة مدرس خصوصي محمول.
\end{itemize}

\booksubsection{النماذج الكبيرة على الجهاز: الهاتف يشغل الذكاء الاصطناعي بنفسه}

أصبحت أحدث الهواتف الرائدة قادرة على تشغيل نماذج ذكاء اصطناعي صغيرة محليًا (عادة بحوالي 7 مليارات معامل)، والتغييرات الناتجة عن ذلك هي:

\begin{itemize}
  \item \textbf{لا حاجة للإنترنت}: استخدام ميزات الذكاء الاصطناعي حتى دون اتصال، في الطائرة أو في المترو.
  \item \textbf{خصوصية أكثر أمانًا}: محادثاتك وبياناتك لا تُرفع إلى السحابة، بل تُعالج بالكامل محليًا على الهاتف.
  \item \textbf{استجابة أسرع}: بدون زمن انتقال للشبكة، تكون ردود الذكاء الاصطناعي فورية تقريبًا.
  \item \textbf{القيود الحالية}: لا تزال قدرة النماذج على الجهاز أقل من النماذج الكبيرة السحابية، والمهام المعقدة لا تزال تتطلب اتصالًا. لكن مع تحسن قوة الحوسبة للذكاء الاصطناعي في الرقائق، ستصبح النماذج على الجهاز أقوى فأقوى.
\end{itemize}

\booksubsection{تكامل الذكاء الاصطناعي على مستوى النظام:}

\begin{itemize}
  \item \textbf{Apple Intelligence من آبل}: نظام ذكاء اصطناعي قُدِّم في iOS 18، يدمج أدوات الكتابة وعصا الصور السحرية وتحسين Siri وملخصات الإشعارات وغيرها. يعطي الأولوية للمعالجة على الجهاز، ولا يلجأ إلى السحابة إلا للمهام المعقدة.
  \item \textbf{نموذج Pangu الكبير من هواوي}: مدمج في نظام HarmonyOS، حصل مساعد Celia على قدرات النموذج الكبير، ويدعم تلخيص المستندات والتعرف الذكي على الشاشة وتحرير الصور بالذكاء الاصطناعي وغيرها.
  \item \textbf{HyperAI من شاومي}: قدرات ذكاء اصطناعي على مستوى النظام، تدعم البحث عن الصور بالذكاء الاصطناعي والملخصات بالذكاء الاصطناعي والترجمة الفورية للعناوين بالذكاء الاصطناعي والحوار مع النموذج الكبير على الجهاز وغيرها.
  \item \textbf{Galaxy AI من سامسونج}: الترجمة الفورية للمكالمات وملخصات الملاحظات بالذكاء الاصطناعي وتوليد الصور وتحريرها والبحث الدائري وغيرها من الوظائف.
\end{itemize}

\booksection{خامسًا: بعض سيناريوهات التطبيق النموذجية والتطبيق العملي}

لكي تشعر بشكل أكثر حدسية بـ “القوة الخارقة” للذكاء الاصطناعي على الهاتف، سنستخدم خمسة سيناريوهات محددة أدناه لنوضح كيف “ينقذ الموقف” الذكاء الاصطناعي في الحياة اليومية:

\booksubsection{السيناريو 1: الترجمة الفورية — مترجمك الشخصي أثناء السفر للخارج}

\textbf{الأداء الفعلي}: في باريس ولا تعرف الفرنسية لطلب الطعام؟ يدعم تطبيق “الترجمة” من iPhone وضع المحادثة الثنائية وجهًا لوجه، حيث يرى كل طرف الترجمة بلغته، بل يمكنه الترجمة الفورية الصوتية. أما في معسكر Android، فإن Galaxy AI من سامسونج يتيح الترجمة الفورية للمكالمات، فتتحدث بالصينية على الهاتف ويسمع الطرف الآخر الإنجليزية، والعكس صحيح.

\textbf{المبدأ الكامن}: مسار من ثلاث خطوات: التعرف على الكلام (ASR) الترجمة الآلية (NMT) تركيب الكلام (TTS)، وقد أصبح زمن الانتقال من البداية إلى النهاية قابلًا للتحكم في أقل من ثانية.

\booksubsection{السيناريو 2: الإزالة بالذكاء الاصطناعي — إنقاذ صور السائقين “التي أفسدها المارة”}

\textbf{الأداء الفعلي}: التقطت صورة ذاتية رائعة في معلم سياحي، لكن الخلفية مزدحمة بالسياح الآخرين؟ استخدم “الإزالة السحرية” لتحديد المارة، وسيتعرف الهاتف تلقائيًا على المارة و“يمحوهم”، ثم يملأ الخلفية بذكاء، وغالبًا ما تكون النتائج مذهلة.

\textbf{المبدأ الكامن}: يحدد نموذج تقسيم الصور (Inpainting) المنطقة المحددة ويحذفها أولًا، ثم “يستنتج” المحتوى المعقول لملئها بناءً على وحدات البكسل المحيطة والفهم الدلالي. حجم بيانات التدريب لنموذج الذكاء الاصطناعي هذا يصل إلى ملايين الصور.

\booksubsection{السيناريو 3: تحرير الصور بالذكاء الاصطناعي — بضغطة واحدة تصبح “أستاذ تصوير”}

\textbf{الأداء الفعلي}: التقطت صورة منظر طبيعي رمادية في يوم غائم؟ يمكن لـ “السماء السحرية” بالذكاء الاصطناعي استبدال السماء بضغطة واحدة بسماء صافية أو غروب أو حتى سماء مرصعة بالنجوم. ويمكن لتوسيع الصور بالذكاء الاصطناعي تحويل التكوين العمودي إلى تكوين أفقي، حيث “يستنتج” الهاتف المشهد خارج الإطار.

\textbf{المبدأ الكامن}: نموذج الانتشار (Diffusion Model) القائم على الذكاء الاصطناعي التوليدي (Generative AI)، يفهم المحتوى الدلالي للصورة ثم يولد وحدات بكسل جديدة. هذا هو “الحد الأقصى” الحالي لمعالجة الصور بالذكاء الاصطناعي.

\booksubsection{السيناريو 4: معالجة المستندات — مضاعف الكفاءة للموظفين}

\textbf{الأداء الفعلي}: بعد انتهاء الاجتماع، ألقِ تسجيل الاجتماع إلى الذكاء الاصطناعي، فيقوم تلقائيًا بتنظيم محاضر اجتماع منظمة هيكليًا واستخراج المهام وتوليد ملخص. تقرير PDF إنجليزي من 10 صفحات، يمكن للذكاء الاصطناعي استخلاص الاستنتاجات الأساسية لك في 10 ثوانٍ.

\textbf{المبدأ الكامن}: النماذج الكبيرة ذات السياق الطويل (Long Context LLM) تستطيع “قراءة” مستندات بعشرات الآلاف من الكلمات وإجراء استدلال متعدد الجولات واستخراج المعلومات. يتمتع كل من Apple Intelligence وCelia من هواوي وWPS AI وYuanbao من Tencent وغيرها بهذه القدرة.

\booksubsection{السيناريو 5: البحث عن الصور والتعرف على الأشياء بالذكاء الاصطناعي — “محرك البحث البصري” الخاص بك}

\textbf{الأداء الفعلي}: رأيت زهرة لا تعرفها في الشارع، افتح الكاميرا ليخبرك الذكاء الاصطناعي فورًا “هذه زهرة الكركديه”. تبحث في ألبوم الصور عن صور التقطتها العام الماضي على شاطئ البحر، ابحث عن “البحر”، فيجد الذكاء الاصطناعي تلقائيًا جميع الصور التي تحتوي على البحر دون تصفح يدوي.

\textbf{المبدأ الكامن}: تقوم نماذج الرؤية الحاسوبية (CV) بتحويل محتوى الصور إلى متجهات (Embedding)، ثم تطابقها مع الاستعلامات النصية من حيث التشابه. هذا “البحث متعدد الوسائط” هو قفزة كبيرة في إدارة ألبوم الصور على الهاتف.

\booksection{سادسًا: تحسين المطالبات — “صيغة التواصل” التي تجعل الذكاء الاصطناعي يطيعك}

تعتمد جودة إجابات الذكاء الاصطناعي إلى حد كبير على كيفية “طرحك للسؤال”. نفس السؤال، بطرق مختلفة، يعطي إجابات متباينة تمامًا. سنقدم أدناه\textbf{نهج تواصل “الدور-الخلفية-المهمة-القيود” العملي والسهل الاستخدام}، مع الكثير من الأمثلة.

\begin{tcbinfo}{لماذا نحتاج إلى إطار عمل؟}
كثير من الناس يستخدمون الذكاء الاصطناعي مثل استخدام Baidu — يكتبون بضع كلمات مفتاحية. لكن الذكاء الاصطناعي ليس محرك بحث، بل هو نموذج لغوي “يجيب بعد فهم المقروء”.\textbf{إذا أعطيته سياقًا أكثر ووضحت ما تريده ومنعته عن الانحراف}، فسيعطيك نتائج دقيقة. الإطار هو ترجمة “كيف تبدو الإجابة الجيدة” في ذهنك إلى تعليمات يفهمها الذكاء الاصطناعي.

\end{tcbinfo}

\booksubsection{الدور (Role): أخبر الذكاء الاصطناعي من هو}

أعطِ الذكاء الاصطناعي “هوية” محددة، فيستدعي المعرفة المهنية وأسلوب التعبير في المجال المقابل.\textbf{كلما كان تحديد الدور أكثر تحديدًا، كان المخرج أكثر احترافية}.

\booksubsection{أمثلة عملية:}

\begin{itemize}
  \item "أنت\textbf{مهندس صيانة هواتف شاومي بخبرة 10 سنوات}." (مشكلة تقنية)
  \item "أنت\textbf{معلم فيزياء في المرحلة الإعدادية متخصص في تبسيط المفاهيم المعقدة}." (مشكلة علمية)
  \item "أنت\textbf{محرر كتب ذو ذوق جمالي وقراءات كثيرة وأسلوبه الخاص}." (مشكلة كتابية)
  \item "أنت\textbf{مدون طبخ يعد الإفطار يوميًا للموظفين}." (مشكلة حياتية)
  \item "أنت\textbf{طبيب صيني تقليدي صبور ولطيف لا يسخر أبدًا من المبتدئين}." (مشكلة صحية)
\end{itemize}

\textbf{لماذا هو فعال؟}يحدد الدور\textbf{نطاق المعرفة}و\textbf{أسلوب التعبير}و\textbf{العمق المهني}للذكاء الاصطناعي. الذكاء الاصطناعي بلا دور هو “شامل لكنه لا يتقن شيئًا”؛ وبوجود الدور، يستطيع “الدخول” في تلك الهوية والتفكير مثل ذلك المحترف.

\booksubsection{الخلفية (Background): أخبر الذكاء الاصطناعي عن وضعك}

وضّح مسبقًا “من أنت، وما وضعك الحالي، وما قيودك”، فيستطيع الذكاء الاصطناعي تقديم إجابات\textbf{تناسبك فعلًا}بدلًا من الكلام العام.

\booksubsection{أمثلة عملية:}

\begin{itemize}
  \item "هاتفي\textbf{شاومي 14، اشتريته قبل نصف عام، ويسخن بشدة عند لعب الألعاب مؤخرًا}." (حالة الجهاز)
  \item "أنا\textbf{امرأة متقاعدة عمري 60 عامًا، بصري ضعيف نسبيًا، ولا أحسن استخدام الهواتف الذكية}." (ملف المستخدم)
  \item "أنا\textbf{أستعد لامتحان CET-6 في يونيو 2026، ومستواي متوسط}." (الخلفية التعليمية)
  \item "أدير\textbf{مقهى صغيرًا في مدينة من الدرجة الثالثة، إيراده الشهري 30 ألف يوان، وأوظف عاملين}." (الخلفية التجارية)
  \item "أنا\textbf{أنجبت طفلًا للتو وأقوم بالنفاس في المنزل، أعاني من نقص حاد في النوم}." (مرحلة حياتية)
\end{itemize}

\textbf{لماذا هو فعال؟}تتيح معلومات الخلفية للذكاء الاصطناعي معرفة\textbf{“لمن يخدم”}و\textbf{“في أي ظروف”}. نفس السؤال، يُشرح بشكل مختلف للطالب عنه للوالد؛ وتختلف النصيحة بين سكان المدن وسكان البلدات. بدون خلفية، لا يستطيع الذكاء الاصطناعي إلا تقديم إجابات “بالقاسم المشترك الأكبر”، وغالبًا ما تكون غير مؤثرة.

\booksubsection{المهمة (Task): أخبر الذكاء الاصطناعي بوضوح ماذا يفعل}

في هذه الخطوة استخدم\textbf{أفعالًا في البداية}(اكتب، اذكر، قارن، ترجم، اشرح، لخّص...)، وصف المهمة بشكل\textbf{محدد وقابل للتنفيذ}. المهام الغامضة تجلب فقط إجابات غامضة.

\booksubsection{أمثلة عملية:}

\begin{itemize}
  \item "\textbf{اذكر 5 أذونات يجب أن أغلقها فورًا}، واشرح السبب." (إجراء واضح)
  \item "\textbf{ترجم رسالة البريد الإلكتروني الإنجليزية التالية إلى الصينية}، مع الحفاظ على النبرة التجارية." (مهمة محددة)
  \item "\textbf{قارن بين هذه الهواتف الأربعة في فئة 3000 يوان}، وقدم جدول مقارنة للمزايا والعيوب." (تنسيق واضح)
  \item "\textbf{اشرح 'ما هو 5G' بكلمات يفهمها تلاميذ المدرسة الابتدائية}." (أسلوب واضح)
  \item "\textbf{اكتب ثلاث نصوص قصيرة لفيديو لمتجر الإفطار الخاص بي}، تستهدف الموظفين والأزواج والآباء ذوي الأطفال على التوالي." (تقسيم واضح)
\end{itemize}

\textbf{لماذا هو فعال؟}جزء المهمة هو “مركز القيادة” لعمل الذكاء الاصطناعي.\textbf{كلما كانت الأفعال أوضح والطلبات أكثر تحديدًا، عرف الذكاء الاصطناعي “ماذا يسلِّم”}. الكلمات الغامضة مثل “افعل شيئًا” و“جهّز” و“انظر” هي من المحظورات الكبرى مع الذكاء الاصطناعي — فسيفسر نيتك وغالبًا ما ينحرف عن المسار.

\booksubsection{القيود (Constraints): أخبر الذكاء الاصطناعي ما لا يجب أن يفعله}

يبدو “القيود” وكأنه “تقييد”، لكنه في الواقع أداة قوية\textbf{لتقريب الذكاء الاصطناعي من المخرج الأنسب لك}. تشمل الأبعاد الشائعة: عدد الكلمات، التنسيق، النبرة، عناصر التجنب، العناصر الضرورية.

\booksubsection{أمثلة عملية:}

\begin{itemize}
  \item "\textbf{لا تتجاوز 200 كلمة، اذكرها في نقاط}." (الطول + التنسيق)
  \item "\textbf{النبرة يجب أن تكون ودودة وطبيعية، كمحادثة بين الأصدقاء}، لا تستخدم 'أولًا، ثانيًا، ثالثًا' مثل هذه اللغة الرسمية." (النبرة)
  \item "\textbf{لا تستخدم المصطلحات الفنية}، وعند ذكر مصطلح مهني اشرحه عرضًا." (قيود المفردات)
  \item "\textbf{لا تتجاوز كل فقرة 3 أسطر}، لتسهيل القراءة على الهاتف." (التنسيق)
  \item "\textbf{يجب أن تتضمن: سيناريو الاستخدام، المزايا والعيوب، مدى مطابقتها لميزانيتي}." (العناصر الضرورية)
  \item "\textbf{لا تقدم نصائح استثمارية محددة}، اشرح المبادئ فقط." (الحدود)
\end{itemize}

\textbf{لماذا هو فعال؟}بدون قيود، يميل الذكاء الاصطناعي إلى “التحرر الإبداعي” — إما أن يكون أكاديميًا للغاية أو طويلًا للغاية أو خارج الموضوع. القيود الواضحة = “معايير التسليم” الواضحة، ومواءمة التوقعات بين الطرفين تجعل النتائج قابلة للتحكم.

\booksubsection{تقييم جودة المطالبة: 5 خصائص للمطالبة الجيدة}

كتابة مطالبة جيدة ليست أمرًا لمرة واحدة،\textbf{بعد الكتابة تحقق من هذه النقاط الخمس أولًا}قبل إرسالها للذكاء الاصطناعي:

\begin{enumerate}
  \item \textbf{محدد لا مجرد}: ليس "اكتب مقالًا"، بل "اكتب دليل شراء هاتف من 800 كلمة و3 أقسام فرعية موجهاً للأمهات المبتدئات".
  \item \textbf{يتضمن الخلفية لا منعزل}: أخبر الذكاء الاصطناعي "من أنا، ولماذا أسأل هذا"، ليتمكن من تقديم إجابة "مخصصة لك".
  \item \textbf{تنسيق مخرج واضح}: هل هو جدول أم قائمة أم فقرات أم JSON؟ حدد مسبقًا.
  \item \textbf{قيود كافية}: عدد الكلمات، الأسلوب، عناصر التجنب، اختيار واحد من الثلاثة يحقق تحسنًا واضحًا.
  \item \textbf{قابل للتحقق من قبل الذكاء الاصطناعي}: بعد السؤال، اطلب من الذكاء الاصطناعي إعادة سرد طلبك؛ إذا لم يستطع توضيحه فأنت لم توضحه.
\end{enumerate}

\booksubsection{المشكلات الشائعة وحلولها}

\begin{center}
\begin{tabularx}{\linewidth}{|X|X|}
\hline
\textbf{العرض} & \textbf{الحل} \tabularnewline
\hline
الإجابة عامة وغير عملية & أضف خلفية (من أنت/أي سيناريو) \tabularnewline
الإجابة أكاديمية/رسمية للغاية & أضف قيود النبرة + قائمة رفض "اللغة الرسمية" \tabularnewline
الإجابة خارج الموضوع/لا تجيب على السؤال & أعد صياغة المهمة بأفعال، واطلب "أجب عن هذا فقط" \tabularnewline
الإجابة طويلة جدًا لا تكملها & حدد عدد الكلمات بصرامة (أقل من 200 كلمة) \tabularnewline
الذكاء الاصطناعي يتحدث بجدية عن لا شيء & أضف "إذا لم تكن متأكدًا اكتب 'لست متأكدًا'" + تحقق ذاتي من الحقائق المهمة \tabularnewline
النسخة الأولى غير مرضية & لا تعيد الكتابة، ألحق بالمطالبة الأصلية "اجعلها ... أكثر" \tabularnewline
\hline
\end{tabularx}
\end{center}

\booksubsection{ثلاث حالات سيناريو كاملة}

\begin{tcbinfo}{السيناريو الأول: مساعد الحياة (اقتراحات إفطار)}
\textbf{المطالبة}:

أنت مدون إفطار متخصص في التنسيق الغذائي (\textbf{الدور}). أنا أم عاملة، ليس لدي سوى 15 دقيقة فقط كل صباح لإعداد الإفطار، ولدي طفل عمره 6 سنوات وحماتي عمرها 60 عامًا (\textbf{الخلفية}).\textbf{اذكر 7 أيام من إفطار غير مكرر}(\textbf{المهمة}). كل إفطار يجب أن: (1) يستغرق التحضير أقل من 15 دقيقة (2) يكون متوازنًا غذائيًا (بروتين + كربوهيدرات + خضار وفاكهة)(3) مناسبًا للكبار والصغار (4) مكوناته متوفرة في السوبر ماركت العادي (\textbf{القيود}حدد “Add a README file” — README هو “كتيب التعليمات” للمشروع، ويُنصح بوجوده في كل مستودع.

\textbf{لماذا هو فعال}: يحدد الدور المستوى المهني، وتعطي الخلفية القواسم المشتركة للفئات الثلاث، وتحدد المهمة المُسلَّمات بوضوح، وتضمن القيود الأربعة القابلية للتنفيذ.

\end{tcbinfo}

\begin{tcbinfo}{السيناريو الثاني: المساعدة في التعلم (التحضير لامتحان CET-6)}
\textbf{المطالبة}:

أنت\textbf{معلم CET-6 ذو خبرة}(الدور). أنا طالب جامعي عادي،\textbf{حصلت على 480 في CET-4، وهدفي 500 في CET-6، وبقي لدي 3 أشهر}(الخلفية).\textbf{ضع لي خطة تحضير لـ 12 أسبوعًا قابلة للتنفيذ}(المهمة). المتطلبات: (1) توزيع واضح للوقت بين الاستماع والتحدث والقراءة والكتابة أسبوعيًا (2) اتجاه التركيز لكل شهر (3) توصية بـ 3 مراجع مناسبة (مع ذكر السبب بالصينية) (4) الإخراج في شكل جدول (القيود).

\textbf{لماذا هو فعال}: يحدد الدور الخلفية المهنية، وتعطي الخلفية نقاط ارتكاز للنقاط والوقت، وتحدد المهمة “خطة مرحلية” بوضوح، وتتطلب القيود “جدولة” و“قابلية للتنفيذ”.

\end{tcbinfo}

\begin{tcbinfo}{السيناريو الثالث: كفاءة العمل (تنظيم محاضر الاجتماع)}
\textbf{المطالبة}:

أنت\textbf{مساعد إداري كبير}(الدور). انتهيت للتو من اجتماع قسم لمدة ساعة،\textbf{ناقشنا فيه خطة الترويج التسويقي للربع الثاني}(الخلفية).\textbf{نظِّم محاضر الاجتماع التالية في محاضر منظمة هيكليًا}(المهمة). المتطلبات: (1) قسمها إلى ثلاثة أقسام “القرارات / المهام / المخاطر والاعتراضات” (2) حدد المسؤول والموعد النهائي لكل مهمة (إذا لم يُذكر فاكتب “بانتظار التأكيد”) (3) إجمالي الكلمات في حدود 600 كلمة (4) الإخراج بتنسيق markdown (القيود).

\textbf{لماذا هو فعال}: يحدد الدور المنظور الإداري، وتعطي الخلفية موضوع الاجتماع، وتحدد المهمة الإخراج المنظم هيكليًا، وتجعل القيود الأربعة النتيجة قابلة لإعادة الاستخدام مباشرة.

\end{tcbinfo}

\textbf{مستوى متقدم: اجعل الذكاء الاصطناعي “مساعدك الشخصي”}\par

جرب هذه الاستخدامات المتقدمة لتختبر القوة الحقيقية للذكاء الاصطناعي:

\begin{itemize}
  \item \textbf{تحسين الحوار متعدد الجولات}: عند عدم الرضا عن إجابة الذكاء الاصطناعي، يمكنك أن تقول "هذا احترافي للغاية، اشرحه مرة أخرى بكلمات يفهمها تلاميذ المدرسة الابتدائية".
  \item \textbf{سلسلة الأفكار (CoT)}: اجعل الذكاء الاصطناعي "يفكر خطوة بخطوة"، فجودة الإجابات عادة ما تكون أعلى: "حلل هذه المسألة الرياضية خطوة بخطوة".
  \item \textbf{تعاون ذكاءات اصطناعية متعددة}: استخدم A لكتابة المسودة، وB للتحسين، وC للترجمة — التركيبات تعطي نتائج أفضل.
  \item \textbf{المطالبات النظامية}: اضبط مسبقًا هويتك وتفضيلاتك ومحظوراتك في إعدادات الذكاء الاصطناعي، لتكون الإجابات أكثر ملاءمة لاحتياجاتك.
\end{itemize}

\booksection{سابعًا: حدود الذكاء الاصطناعي وتحدياته — حافظ على حكمتك الواعية}

مع أن الذكاء الاصطناعي قوي، إلا أنه ليس “قادرًا على كل شيء”. إن الإدراك الواعي لحدوده يمكننا من استخدامه بشكل أفضل:

\begin{itemize}
  \item \textbf{مشكلة "الهلوسة"}: قد يتحدث الذكاء الاصطناعي بجدية عن لا شيء (المصطلح المهني: Hallucination). قد يختلق أحداثًا تاريخية غير موجودة، أو اقتباسات خاطئة، أو إحصاءات وهمية. يجب التحقق بشريًا من المعلومات المهمة.
  \item \textbf{توقف المعرفة}: لمعظم بيانات تدريب الذكاء الاصطناعي تاريخ انقطاع (مثل أوائل 2024)، فهو لا يعرف ما حدث بعد ذلك.
  \item \textbf{قدرة استدلال محدودة}: الاستدلال متعدد الخطوات المعقد، والإثباتات الرياضية، والفخاخ المنطقية، لا يزال الذكاء الاصطناعي عرضة للخطأ فيها.
  \item \textbf{مخاطر الخصوصية}: في خدمات الذكاء الاصطناعي السحابية، قد تُستخدم محادثاتك في تدريب النماذج. يجب إعطاء الأولوية للذكاء الاصطناعي على الجهاز للمعلومات الحساسة.
  \item \textbf{التحيز والقيم}: قد يحمل الذكاء الاصطناعي تحيزات من بيانات التدريب، لذا يجب الحفاظ على التفكير المستقل.
\end{itemize}

\begin{tcbinfo}{نصائح لاستخدام الذكاء الاصطناعي}
\begin{itemize}
  \item قد يكون المحتوى الذي يولده الذكاء الاصطناعي غير دقيق (مشكلة "الهلوسة")، لذا تحقق دائمًا من المعلومات المهمة.
  \item لا تدخل معلومات خصوصية شخصية إلى الذكاء الاصطناعي (رقم الهوية، كلمات المرور، المعلومات المصرفية، إلخ).
  \item أحسن استخدام المطالبات (Prompt): كلما كان الوصف أكثر تحديدًا ووضوحًا، ارتفعت جودة إجابات الذكاء الاصطناعي.
  \item الذكاء الاصطناعي أداة وليس سلطة، حافظ على تفكيرك المستقل وحكمك.
  \item عند مواجهة قرارات حاسمة (طبية، قانونية، مالية)، استشر محترفين بدلًا من الاعتماد على الذكاء الاصطناعي.
\end{itemize}

\end{tcbinfo}

\begin{tcbpractice}{بكلام بسيط، الأتمتة هي}
\begin{enumerate}
  \item \textbf{جعل الهاتف يؤدي تلقائيًا سلسلة من المهام نيابة عنك عند تحقيق شرط معين.}لا تحتاج إلى التشغيل اليدوي في كل مرة، ولا إلى تذكر ما يجب فعله — فالهاتف يحكم على التوقيت بنفسه وينفذ بنفسه.
  \item \textbf{على سبيل المثال: كل صباح تستيقظ، تغلق يدويًا وضع عدم الإزعاج، وتفتح WeChat لرؤية الرسائل، وتفتح تطبيق الطقس لترى ماذا ترتدي اليوم، وتفتح NetEase Cloud Music لتشغيل أغنية. إذا حزمت هذه الخطوات في تدفق أتمتة واحد، وضبطتها على "التشغيل التلقائي كل صباح الساعة 7"، فعندما تستيقظ ما عليك سوى التقاط الهاتف، وكل شيء يكون جاهزًا. أليس هذا ممتعًا بمجرد التفكير فيه؟}الأتمتة في جوهرها هي سلسلة منطقية "إذا... فـ..."، أي if-else في البرمجة. إذا فهمت طريقة التفكير هذه، حتى لو لم تكتب كودًا طوال حياتك، ستكون أكثر تنظيمًا في حل المشكلات. العناصر الأساسية:
  \item \textbf{شرط التشغيل}: متى يبدأ؟ يمكن أن يكون "كل صباح الساعة 7"، أو "عندما أتصل بشبكة Wi-Fi المنزلية"، أو "عندما أصل إلى مكان ما"، أو "عندما تنخفض بطارية الهاتف عن 20\%".
\end{enumerate}

\end{tcbpractice}

\begin{tcbinfo}{حالة 2: عملية نقل البيانات الكاملة عند تغيير الهاتف}
\begin{itemize}
  \item \textbf{السيناريو}: iPhone 13 Xiaomi 15 تغيير بين علامتين.
  \item \textbf{الخطوات}: (1) نسخ احتياطي للهاتف القديم عبر iCloud/WeChat المدمج (2) تسجيل الخروج من Apple ID + إغلاق "Find" (3) الهاتف الجديد يستخدم تطبيق "Xiaomi Switch" لمسح الرمز (4) اختيار عناصر النقل (5) WeChat يستخدم "نقل سجل الدردشة" المدمج بشكل منفصل (6) التحقق من البيانات الرئيسية (جهات الاتصال/الألبوم/WeChat) (7) مسح مشفّر للهاتف القديم.
  \item \textbf{النتيجة المتوقعة}: فقدان بيانات صفري، الهاتف القديم يمكن بيعه بأمان.
\end{itemize}

\end{tcbinfo}
